%% Chương 1: Mở đầu
\clearpage
\phantomsection
\section*{CHƯƠNG 1. MỞ ĐẦU}
\addcontentsline{toc}{section}{\numberline{}CHƯƠNG 1. MỞ ĐẦU}
\setcounter{section}{1}
\setcounter{subsection}{0}
\setcounter{figure}{0}
\setcounter{table}{0}
\setcounter{equation}{0}

\subsection{Bối cảnh và động lực}

\subsubsection{Sự bùng nổ của video trực tuyến}

Trong thập kỷ vừa qua, lưu lượng video trực tuyến (video streaming) đã tăng trưởng
với tốc độ chưa từng có. Theo các báo cáo thị trường, video hiện chiếm hơn 80\%
tổng lưu lượng Internet toàn cầu, và con số này dự báo sẽ vượt 90\% vào năm 2030.
Sự tăng trưởng này xuất phát từ nhiều xu hướng đồng thời: sự phổ biến của các
nền tảng phát trực tuyến theo yêu cầu (video-on-demand --- VOD) như YouTube, Netflix,
TikTok; sự mở rộng của hội nghị truyền hình (video conferencing) sau đại dịch;
sự phát triển của các hệ thống giám sát an ninh (CCTV, body cam); và sự ra đời
của các ứng dụng thực tế tăng cường (AR) và thực tế ảo (VR).

Trong bức tranh đó, chuẩn nén video H.264/AVC (ITU-T H.264 / ISO/IEC 14496-10)
\cite{itu_h264} giữ vị trí thống trị trong nhiều năm. Mặc dù các chuẩn mới hơn như
H.265/HEVC hay AV1 đang dần được triển khai, H.264 vẫn là lựa chọn mặc định
trên phần lớn thiết bị và nền tảng nhờ sự cân bằng xuất sắc giữa hiệu quả nén,
tốc độ mã hoá/giải mã, và tính tương thích phần cứng. Phần lớn camera giám sát,
thiết bị di động, và hệ thống phát sóng đều xuất ra video ở định dạng H.264.

\subsubsection{Yêu cầu bảo mật nội dung video}

Khi video trở thành phương tiện truyền tải thông tin quan trọng, nhu cầu bảo vệ
nội dung video khỏi truy cập trái phép ngày càng trở nên cấp thiết. Các tình huống
điển hình bao gồm:

\begin{itemize}
  \item \textbf{Bảo vệ bản quyền:} Nội dung giải trí có giá trị cao (phim, thể thao)
        cần được mã hoá để tránh xem lậu (piracy). Các hệ thống DRM (Digital Rights
        Management) như Widevine, FairPlay đều dựa trên mã hoá video stream.

  \item \textbf{Video giám sát an ninh:} Hình ảnh camera tại ngân hàng, sân bay,
        cơ quan chính phủ phải được bảo mật trong quá trình truyền qua mạng để
        tránh rò rỉ thông tin nhạy cảm hoặc bị giả mạo.

  \item \textbf{Hội nghị truyền hình y tế và pháp lý:} Các buổi tư vấn y tế từ xa,
        phiên toà trực tuyến yêu cầu mức độ bảo mật cao theo quy định pháp luật.

  \item \textbf{Truyền dữ liệu trên kênh không tin cậy:} Video truyền qua mạng
        không dây, mạng công cộng có nguy cơ bị nghe lén (eavesdropping)
        hoặc tấn công man-in-the-middle.
\end{itemize}

\subsubsection{Hạn chế của mã hoá toàn bộ luồng bit}

Giải pháp truyền thống là mã hoá toàn bộ luồng bit video bằng AES-128 hoặc AES-256
ở chế độ CBC/CTR. Cách tiếp cận này đảm bảo an toàn tuyệt đối nhưng có những
hạn chế thực tiễn quan trọng:

\begin{itemize}
  \item \textbf{Chi phí tính toán cao:} Video 4K 30fps có thể có bitrate 50--100 Mbps.
        Mã hoá AES toàn bộ bitstream đòi hỏi xử lý $6$--$12$ MB dữ liệu mỗi giây,
        ngốn đáng kể tài nguyên CPU trên thiết bị nhúng không có AES-NI hardware.

  \item \textbf{Độ trễ tăng:} Trong các ứng dụng thời gian thực (live streaming,
        video call), mỗi mili-giây trễ thêm đều ảnh hưởng đến trải nghiệm người dùng.
        Mã hoá toàn bộ có thể tăng độ trễ lên 10--30 ms trên thiết bị nhúng.

  \item \textbf{Không tương thích với codec pipeline:} Một số hệ thống cần
        transcode video ở phía server trong khi vẫn bảo mật nội dung (proxy encryption),
        điều này không khả thi khi toàn bộ payload bị mã hoá.
\end{itemize}

Mã hoá chọn lọc (selective encryption) \cite{cheng_selective} ra đời như một giải pháp
trung gian: chỉ mã hoá phần quan trọng nhất của video (perceptually significant
components), giảm chi phí tính toán mà vẫn đảm bảo mức độ bảo vệ chấp nhận được.

\subsection{Vấn đề nghiên cứu}

\subsubsection{Hệ số DC trong H.264 --- vùng nhạy cảm nhất}

Trong chuẩn H.264/AVC, mỗi khối ảnh (macroblock) được nén qua ba bước chính:
dự đoán nội (intra prediction), biến đổi cosine rời rạc (DCT), và lượng tử hoá
(quantization). Sau quá trình này, hệ số DC (coefficient ở tần số 0---0) của
biến đổi DCT là giá trị thể hiện mức cường độ trung bình của khối ảnh.
Hệ số DC quyết định độ sáng/tối cơ bản của vùng ảnh tương ứng.

Từ góc độ bảo mật tri giác (perceptual security), hệ số DC là mục tiêu lý tưởng
vì hai lý do:
\begin{enumerate}
  \item \textbf{Tác động thị giác lớn:} Nhiễu loạn hệ số DC gây ra distortion
        có thể nhìn thấy ngay cả ở bitrate cao. PSNR của video bị nhiễu loạn DC
        giảm xuống khoảng 7--8 dB, tương đương mức méo hình nghiêm trọng.
        Thử nghiệm của chúng tôi trên file output.h264 xác nhận điều này.

  \item \textbf{Chi phí thấp:} DC bytes chiếm 25 byte/NALU --- dưới 0.01\%
        kích thước bitstream trung bình. Mã hoá chúng tiêu tốn không đáng kể
        so với decode toàn bộ video.
\end{enumerate}

\subsubsection{Điểm yếu của các thuật toán mã hoá stream thông thường}

Khi áp dụng AES-CTR, DES-OFB, hoặc RC4 vào mã hoá chọn lọc DC, một vấn đề thiết
kế quan trọng thường bị bỏ qua: \textbf{keystream reuse}. Nếu cùng khóa K và cùng
IV/counter/seed được dùng cho mọi NALU, thì:

\begin{equation}
\text{Enc}(P_1, K) = P_1 \oplus ks \quad\text{và}\quad \text{Enc}(P_2, K) = P_2 \oplus ks
\end{equation}

Nếu kẻ tấn công biết $P_1$ (DC bytes của một NALU thông qua phân tích thống kê),
hắn có thể tính $ks = \text{Enc}(P_1, K) \oplus P_1$, từ đó giải mã mọi NALU khác.

Vấn đề thứ hai là thiếu \textbf{diffusion}: nếu 1 bit của plaintext thay đổi,
chỉ 1 bit ciphertext thay đổi theo (không có hiệu ứng lan truyền). Điều này dẫn
đến NPCR (Number of Pixels Change Rate) thấp --- chỉ $\approx 4\%$ (1/25 bytes),
trong khi mã hoá mạnh cần NPCR $> 99.6\%$.

\subsubsection{Khoảng trống nghiên cứu}

Qua khảo sát các công trình liên quan (phần~\ref{subsec:related_work}), chúng tôi
nhận thấy phần lớn nghiên cứu mã hoá chọn lọc H.264 tập trung vào:
(a) lựa chọn vùng mã hoá (entropy coding, motion vectors, DC/AC),
(b) duy trì tính stream-compliant,
nhưng ít chú ý đến \textbf{per-NALU keystream variation} và \textbf{intra-NALU diffusion}
--- hai tính chất cần thiết để đảm bảo plaintext sensitivity (NPCR/UACI).

Nghiên cứu này lấp đầy khoảng trống đó bằng cách đề xuất thuật toán hybrid
kết hợp ánh xạ hỗn độn (PLCM) tạo keystream khác mỗi NALU, hoán vị Arnold,
và cơ chế khuếch tán chuỗi (chained diffusion feedback).

\subsection{Công trình liên quan}
\label{subsec:related_work}

\subsubsection{Mã hoá chọn lọc video H.264}

Cheng và cộng sự \cite{cheng_selective} đề xuất phân loại mã hoá chọn lọc
thành 3 cấp độ: header encryption, DC coefficient encryption, và motion vector
encryption. Nghiên cứu chứng minh rằng chỉ cần mã hoá DC là đủ để làm mờ nội dung
với chi phí dưới 5\% tổng chi phí mã hoá.

Shahid và cộng sự \cite{shahid_selective} nghiên cứu mã hoá chọn lọc cho
H.264 với yêu cầu duy trì định dạng (format compliance) --- video sau mã hoá
vẫn có thể được xử lý bởi các thiết bị H.264 tiêu chuẩn. Đây là ràng buộc
chặt hơn so với cách tiếp cận của chúng tôi.

\subsubsection{Mã hoá video dựa trên lý thuyết hỗn độn}

Chen và cộng sự \cite{chen_chaos_video} đề xuất sử dụng kết hợp logistic map
và Arnold Cat Map cho mã hoá toàn bộ frame video. Nghiên cứu chứng minh rằng
hệ thống hỗn độn kết hợp đạt NPCR $> 99.6\%$ và Entropy $\approx 8.0$.
Tuy nhiên, nghiên cứu đó áp dụng cho ảnh tĩnh, không tối ưu cho bitstream H.264.

Nghiên cứu của chúng tôi kế thừa ý tưởng kết hợp PLCM và Arnold từ \cite{chen_chaos_video}
nhưng thích nghi cho cấu trúc NALU của H.264 với cơ chế per-NALU nalu\_index
--- điểm khác biệt then chốt.

\subsubsection{NPCR/UACI như chỉ số đánh giá}

Wu và cộng sự \cite{wu_npcr} hệ thống hoá các chỉ số NPCR và UACI, xác lập
ngưỡng lý thuyết cho mã hoá tốt: NPCR $> 99.6\%$ và UACI $\approx 33.46\%$.
Đây là framework đánh giá được chúng tôi áp dụng trực tiếp.

\subsection{Mục tiêu nghiên cứu}

Nghiên cứu này nhằm đạt các mục tiêu cụ thể sau:

\begin{enumerate}
  \item \textbf{Đề xuất thuật toán hybrid:} Thiết kế và cài đặt thuật toán
        mã hoá chọn lọc DC H.264 kết hợp PLCM, Arnold 2D Cat Map, và
        cơ chế khuếch tán chuỗi (chained diffusion feedback), đảm bảo
        per-NALU keystream variation.

  \item \textbf{So sánh hệ thống:} Cài đặt ba thuật toán so sánh (AES-128-CTR,
        DES-OFB, RC4) trong cùng pipeline, đánh giá trên các chỉ số đồng nhất:
        bảo mật (NPCR, UACI, Entropy, SSIM, TSSIM), hiệu năng (ns/NALU), và
        tính đúng đắn (BER, SHA-256).

  \item \textbf{Phân tích hai chiến lược:} Đánh giá S1 (toàn bộ VCL) và S2
        (chỉ I-frame) về hiệu năng và bảo mật, làm rõ trade-off giữa phạm vi
        mã hoá và chi phí tính toán.

  \item \textbf{Thiết kế kiến trúc mở rộng:} Xây dựng IEncryption interface
        (C++) cho phép tích hợp thuật toán mới mà không thay đổi pipeline
        --- hướng tới triển khai FPGA trong tương lai.
\end{enumerate}

\subsection{Phạm vi nghiên cứu}

Để đảm bảo tính tập trung và khả năng tái hiện, nghiên cứu này được giới hạn
trong phạm vi sau:

\begin{itemize}
  \item \textbf{Định dạng video:} H.264 raw bitstream (Annex B format). Không
        hỗ trợ container (MP4, MKV) hoặc các chuẩn khác (H.265, AV1, VP9).

  \item \textbf{Vùng mã hoá:} Chỉ DC bytes tại offset 19, kích thước 25 bytes
        trong mỗi NALU VCL. Không mã hoá SPS, PPS, SEI, motion vectors,
        AC coefficients, hay Cabac data.

  \item \textbf{Dataset:} Một file H.264 duy nhất (\texttt{output.h264}),
        kích thước 219 MB, 19.348 NALU tổng cộng. Đây là hạn chế được thừa
        nhận (xem Chương 7).

  \item \textbf{Môi trường thực thi:} Linux Ubuntu 22.04 (x86-64), GCC 11,
        Python 3.12. Không bao gồm FPGA triển khai phần cứng (là hướng phát triển).

  \item \textbf{Mô hình đe doạ:} Nghiên cứu tập trung vào \emph{perceptual security}
        (làm mờ nội dung thị giác) và \emph{plaintext sensitivity}.
        Không phân tích IND-CCA, side-channel attacks, hay key management.
\end{itemize}

\subsection{Đóng góp của nghiên cứu}

Các đóng góp chính của nghiên cứu này bao gồm:

\begin{enumerate}
  \item \textbf{Thuật toán hybrid mới:} Đề xuất kết hợp PLCM + Arnold 2D + XOR
        + chained diffusion vào mã hoá DC H.264, đạt NPCR $> 99.69\%$ (S2)
        và Entropy $\approx 8.0$ --- vượt ngưỡng tiêu chuẩn.

  \item \textbf{Phân tích pipeline paradox:} Lần đầu tiên mô tả và giải thích
        hiện tượng ``pipeline convergence'' trong S2 --- khi DC $< 1\%$
        tổng thời gian pipeline, tốc độ thuật toán không ảnh hưởng đến thời gian tổng.

  \item \textbf{Bộ dữ liệu thực nghiệm đầy đủ:} Cung cấp kết quả 12 chỉ số
        cho 8 tổ hợp (4 thuật toán $\times$ 2 chiến lược) trên cùng một file
        H.264, cho phép so sánh công bằng (apples-to-apples comparison).

  \item \textbf{Kiến trúc phần mềm mở rộng:} IEncryption interface cho phép
        thêm thuật toán mới với 0 thay đổi pipeline --- nền tảng cho triển khai
        FPGA và nghiên cứu tiếp theo.
\end{enumerate}

\subsection{Cấu trúc báo cáo}

Phần còn lại của báo cáo được tổ chức như sau:

\begin{itemize}
  \item \textbf{Chương 2 --- Cơ sở lý thuyết:} Trình bày nền tảng kiến thức về
        chuẩn H.264/AVC (cấu trúc NALU, hệ số DC, cơ chế nén intra), các nguyên
        lý mật mã hỗn độn (PLCM, Arnold Cat Map), các thuật toán mã hoá cổ điển
        (AES, DES, RC4), và framework đánh giá (NPCR, UACI, Entropy, SSIM, BER).

  \item \textbf{Chương 3 --- Phương pháp đề xuất:} Mô tả chi tiết thiết kế
        thuật toán hybrid, bao gồm PLCM keystream generation, Arnold permutation,
        XOR và chained diffusion, giải mã, phân tích bảo mật và keyspace.

  \item \textbf{Chương 4 --- Các thuật toán so sánh:} Phân tích AES-128-CTR,
        DES-OFB, RC4 từ góc độ điểm yếu bảo mật trong ngữ cảnh mã hoá chọn lọc,
        giải thích lý do NPCR/UACI không áp dụng cho các thuật toán này.

  \item \textbf{Chương 5 --- Thiết lập thực nghiệm:} Mô tả môi trường phần
        cứng/phần mềm, dataset, pipeline C++ end-to-end, và quy trình đo metrics.

  \item \textbf{Chương 6 --- Kết quả và thảo luận:} Phân tích toàn bộ kết quả
        thực nghiệm về tính đúng đắn, hiệu năng, và bảo mật cho tất cả
        8 tổ hợp thuật toán-chiến lược.

  \item \textbf{Chương 7 --- Kết luận:} Tổng kết đóng góp, nhận xét hạn chế
        của nghiên cứu, và đề xuất hướng phát triển trong tương lai.
\end{itemize}
